\documentclass[12pt,a4paper]{article}
% ============================================================
% Preámbulo común a todos los documentos del Proyecto 2
% ============================================================
\usepackage[utf8]{inputenc}
\usepackage[T1]{fontenc}
\usepackage[spanish]{babel}
\usepackage[margin=2.3cm]{geometry}
\usepackage{amsmath}
\usepackage{booktabs}
\usepackage{float}
\usepackage{array}
\usepackage{enumitem}
\usepackage{tikz}
\usetikzlibrary{shapes.geometric, arrows.meta, positioning}
\usepackage[hidelinks]{hyperref}

\renewcommand{\contentsname}{Índice}
\renewcommand{\figurename}{Figura}
\renewcommand{\tablename}{Tabla}

\setlist{itemsep=2pt, topsep=4pt}

\tikzset{
  bloque/.style={rectangle, rounded corners, draw=black, thick, fill=blue!8,
                 text width=7.5cm, align=center, font=\small, inner sep=6pt},
  bloqueMCU/.style={bloque, fill=green!12},
  bloqueSPI/.style={bloque, fill=orange!25},
  bloqueSBC/.style={bloque, fill=violet!12},
  flujo/.style={-{Stealth[length=2.5mm]}, thick},
  bidir/.style={{Stealth[length=2.5mm]}-{Stealth[length=2.5mm]}, thick},
  inicio/.style={rectangle, rounded corners=6pt, draw=black, thick, fill=green!15,
                 minimum width=3.5cm, minimum height=0.8cm, align=center, font=\small},
  proc/.style={rectangle, draw=black, thick, fill=blue!8, text width=5.2cm,
               align=center, font=\small, minimum height=0.9cm, inner sep=4pt},
  dec/.style={diamond, aspect=2.4, draw=black, thick, fill=yellow!20,
              text width=3.0cm, align=center, font=\footnotesize, inner sep=1pt},
}

% Datos del proyecto (un solo lugar para mantenerlos)
\newcommand{\Universidad}{Universidad Nacional de Colombia --- Sede Manizales}
\newcommand{\Facultad}{Facultad de Ingeniería y Arquitectura}
\newcommand{\Carrera}{Carrera de Ingeniería Electrónica}
\newcommand{\Materia}{Programación de sistemas Linux embebido}
\newcommand{\TituloProyecto}{Sistema Híbrido de Adquisición Analógica mediante SPI\\[0.3cm]y Dashboard Embebido en Linux}
\newcommand{\FechaDoc}{Manizales, noviembre de 2026}

% Comando de color-neutral para valores a completar por el equipo
\newcommand{\completar}[1]{\textit{[#1]}}
% ============================================================
% Portada común: \portada{Título del documento}
% ============================================================
\newcommand{\portada}[1]{%
\begin{titlepage}
\centering
{\large \Universidad\par}
{\large \Facultad\par}
{\large \Carrera\par}
\vspace{1.2cm}
{\large Materia: \textbf{\Materia}\par}
\vspace{1.8cm}
{\LARGE\bfseries \TituloProyecto\par}
\vspace{0.5cm}
{\Large #1\par}
\vspace{1.5cm}

\small
\begin{tabular}{@{}p{0.5cm} p{4.2cm} p{2.8cm} p{6.0cm}@{}}
\toprule
& \textbf{Integrante} & \textbf{C.C.} & \textbf{Rol técnico} \\
\midrule
1. & \hrulefill & \hrulefill & Ingeniero de Documentación Técnica y Gestión de Configuración \\[0.7cm]
2. & \hrulefill & \hrulefill & Arquitecto de Sistemas Embebidos \\[0.7cm]
3. & \hrulefill & \hrulefill & Ingeniero de Software y Firmware Embebido \\[0.7cm]
4. & \hrulefill & \hrulefill & Ingeniero de Verificación y Validación (V\&V) \\[0.7cm]
\bottomrule
\end{tabular}

\vfill
\FechaDoc
\end{titlepage}
}

\begin{document}

\portada{Documento 6: Registro de conversaciones con la IA}

\tableofcontents
\newpage

\section{Convención y plantilla}

Herramienta: Claude Code. Modelo: Claude Sonnet 5.5. Los mensajes del usuario se transcriben \textbf{literalmente}; solo se retiraron las credenciales (indicado como \texttt{[omitido]}). Las respuestas del asistente se transcriben literalmente cuando se indica «literal»; el trabajo con herramientas (lectura de archivos, comandos, compilación) se resume en tablas. La transcripción completa generada por la herramienta se obtiene con \texttt{exportar\_chats.ps1}.

\paragraph{Plantilla para nuevas sesiones.}

\begin{quote}\small
\textbf{Sesión N --- fecha --- objetivo}\\
\textbf{Usuario (literal):} \ldots\\
\textbf{Respuesta del asistente (literal):} \ldots\\
\textbf{Acciones con herramientas (resumen):} \ldots\\
\textbf{Verificación humana realizada:} \ldots\\
\textbf{Resultado aceptado o rechazado y por qué:} \ldots
\end{quote}

\section{Sesión 1: revisión, diseño, firmware, software y repositorio}

\subsection{Turno 1 (1 de noviembre de 2026)}

\paragraph{Usuario (literal).}
\begin{quote}
Quiero que leas la documentacion que hay dentro de la carpeta y me digas que informacion adicional necesitas
\end{quote}

\paragraph{Acciones con herramientas.} Lectura de los dos documentos \LaTeX{} originales y del enunciado del profesor.

\paragraph{Respuesta del asistente (literal en los puntos 1 a 13; abreviada en los puntos 14 a 17 y sin la sección final de observaciones menores).}
\begin{quote}\small
Leí los tres archivos: la especificación técnica, la matriz de verificación y el enunciado del profesor. El diseño está bastante completo en papel. Para pasar a implementación o a cerrar los documentos me falta lo siguiente.

\textbf{Hardware y datos concretos.}
\begin{enumerate}[leftmargin=*]
\item \textbf{Raspberry Pi:} el documento dice «Raspberry Pi Model B». ¿Cuál es exactamente (B original, 3B, 3B+, 4B)? Cambia el SoC, la velocidad real de SPI y la versión de Raspberry Pi OS.
\item \textbf{ESP32-S3:} ¿qué placa tienen (DevKitC-1, otra) y qué framework usan (ESP-IDF o Arduino)? Para 40 kSa/s con doble buffer y SPI esclavo conviene ESP-IDF.
\item \textbf{Señales de entrada:} ¿qué rango de tensión y qué frecuencias máximas tienen? ¿Hay acondicionamiento (divisor, offset, protección) o entran directo al ADC de 0--3,3 V?
\item \textbf{Resolución del ADC:} ¿12 bits alineados a 16? ¿Tienen previsto calibrar la no linealidad del ADC del ESP32-S3?
\end{enumerate}

\textbf{Decisiones técnicas abiertas o inconsistentes.}
\begin{enumerate}[leftmargin=*, start=5]
\item \textbf{Protocolo de comandos SPI:} no hay formato definido para el cambio de preset, para la confirmación («ACK») ni para solicitar una ventana. Tampoco se dice cómo sabe la Raspberry Pi que hay una ventana lista, ya que el esclavo no puede iniciar. ¿Quieren una línea GPIO de interrupción (DATA\_READY) o sondeo?
\item \textbf{Formato de la trama:} el encabezado es de 6 bytes (ID, secuencia, longitud y CRC-16), pero no están definidos los campos ni su orden. ¿Lo defino yo?
\item \textbf{Dos canales «simultáneos»:} el ADC del ESP32-S3 multiplexa canales, así que no son realmente simultáneos. ¿Aceptan un desfase entre CH1 y CH2 o necesitan muestreo simultáneo?
\item \textbf{Posible incoherencia en la meta de 40 kSa/s:} el ADC1 del ESP32-S3 en modo continuo (DMA) llega a $\sim$83 kSa/s por canal, pero el modo oneshot es mucho más lento. Hay que fijar el modo de adquisición.
\item \textbf{Refresco vs. captura:} la Raspberry Pi pide una ventana cada 1--4 s, mientras que el ESP32 genera ventanas a 6,4 ms--640 ms. Entonces se descartan casi todas. ¿Es lo esperado o quieren captura bajo demanda (disparo al pedir)? Esto afecta RNF-02, RNF-06 y la prueba IT-07 («sin pérdida de ventanas»).
\item \textbf{Cableado y nivel SPI:} ¿cuáles son los pines concretos? El \texttt{spidev} de la Raspberry Pi a 20 MHz con un ESP32-S3 esclavo suele ser inestable. Su documento da 10--20 MHz como rango, pero habrá que validar con cables cortos.
\end{enumerate}

\textbf{Software y entorno.}
\begin{enumerate}[leftmargin=*, start=11]
\item \textbf{Servidor web:} ¿Nginx o Apache? ¿Prefieren una librería de gráficas concreta (Chart.js, uPlot, canvas puro) o puedo elegirla?
\item \textbf{Lenguaje del servicio en la Raspberry Pi:} ¿Python (spidev) o C? Para 40 kSa/s Python probablemente alcance, pero hay que medirlo.
\item \textbf{Analizador lógico / osciloscopio:} el plan de pruebas lo exige. ¿Con qué instrumentos cuentan para verificar?
\end{enumerate}

\textbf{Entregables y organización.}
\begin{enumerate}[leftmargin=*, start=14]
\item \textbf{Equipo:} la portada deja 4 integrantes y el enunciado permite hasta 6. ¿Cuántos son y qué rol tiene cada uno?
\item \textbf{Repositorio:} ¿ya existe? ¿Dónde se aloja (GitHub, GitLab) y qué estructura quieren?
\item \textbf{Declaración de uso de IA:} requiere los chats completos. ¿Quieren que prepare una plantilla y un registro de esta conversación?
\item \textbf{Estado del trabajo:} ¿qué falta exactamente que haga yo? Por ejemplo código del firmware, servicio de la Raspberry Pi, dashboard, pulir los \texttt{.tex}, o rellenar la matriz con resultados reales.
\end{enumerate}

Si me responden lo más crítico (1, 2, 5, 9, 11, 12 y 17), puedo empezar de inmediato. Lo demás lo asumo con valores razonables y lo dejo anotado.
\end{quote}

\paragraph{Corrección posterior.} El punto 8 decía «$\sim$83 kSa/s por canal». El valor correcto es $\sim$83 mil conversiones por segundo \emph{en total}, es decir $\approx$ 41,7 kSa/s por canal con dos canales. Se corrigió en el diseño (Documento 1).

\subsection{Turno 2 (1 y 2 de noviembre de 2026)}

\paragraph{Usuario (literal).}
\begin{quote}
La raspberry que estamos usando es la 4 model B, la esp32-s3 es la N16R8, en base a estos dos ten en cuenta su rango de frecuencias maximas y resolucion del ADC, el protocolo SPI va entre la raspberry pi y la esp32-s3 este protocolo le va a pasar los datos que se midan en la esp32-s3 a la raspberry, quiero que definas el formato de la trama, respecto al ADC del esp32-s3 que no son realmente simultaneos elije tu la mejor opcion que tenemos, fija el modo de adquisicion, queriamos que fuera con refresco de ventanas, no con captura de bajo demanda, las conexiones que quisiera hacer seria dejar dos jumpers para el CH1(GPIO1) y CH2(GPIO2), y entre la esp32-s3 y la raspberry pi, SCK(GPIO12)->SCLK(GPIO11), MOSI(GPIO11)->MOSI(GPIO10), MISO(GPIO13)->MISO(GPIO9), CS(GPIO10)->CE0(GPIO8), GND(GND)->GND, quiero que tu decidas que libreria de graficas usar para el servidor web, ademas ten en cuenta que tambien los datos se deben visualizar en el monitor serial, para programar la raspberry si C alcanza los 40kSa/s entonces elije C, si no elije python, dame una recomendacion para hacer las pruebas, somos 4 integrantes uno se encargo de la documentacion, otro de la arquitectura, otro de la programacion, otro de las pruebas, a estos colocale un nombre mas tecnico, Ya tengo un repositorio en mi escritorio listo para usarse para subir los cambios se subiran a (\url{https://github.com/Juaromero/Programacion_de_Sistemas_Linux_Embebidos}) ademas quiero que lo organices en 3 subcarpetas que va a tener Proyecto1 (Este despues te lo doy), Proyecto2 (Osciloscopio), Proyecto3 (Este despues te lo doy), mira como lo hice de ejemplo con este otro repositorio mio, hazlo asi (\url{https://github.com/Juaromero/Sistemas_En_Tiempo_Real}), ahora con los chats registralos completos prepara una plantilla y un registro de conversacion, ahora la esp32-s3 la vamos a programar por zephr, ya tengo el entorno de linux WSL conectado a Zephyr y ya esta lista para usarse y ademas la raspberry pi esta conectada a mi PC desde SSH, quiero que me ayudes con la documentacion, programacion y testeo, ademas ayudame subiendo las cosas al github
\end{quote}

\paragraph{Acciones con herramientas (resumen).}

\begin{table}[H]
\centering
\small
\begin{tabular}{@{}p{0.6cm} p{7.2cm} p{7.2cm}@{}}
\toprule
\textbf{\#} & \textbf{Acción} & \textbf{Resultado} \\
\midrule
1 & Lectura del repositorio de ejemplo (subagente de lectura web) & Estructura \texttt{Codigo/}, \texttt{Documentos/}, \texttt{Imagenes/} y README adoptada \\
2 & Exploración del equipo: repositorio local, WSL con Zephyr 4.5.99 y SDK 1.0.1, \texttt{usbipd} & Repositorio local vacío y remoto sin commits \\
3 & Lectura del código fuente de Zephyr (SPI esclavo, ADC con DMA, pinctrl, pruebas de referencia) & Esclavo SPI por FIFO (64 B por CS); ADC con DMA entre 611 y 83\,333 conv/s; pines IOMUX de SPI2 \\
4 & Diseño del protocolo de 5 segmentos de 64 bytes y de 4 presets & Documento 3 \\
5 & Escritura y compilación del firmware & Compila sin errores; 202 KB de flash \\
6 & \texttt{west flash} se bloqueó; diagnóstico con \texttt{usbipd list} & El CH343 estaba compartido pero no adjunto a WSL; se adjuntó con \texttt{usbipd attach}; la consola pasó a UART0 \\
7 & Carga y prueba de los 4 presets por monitor serie & Duraciones de ventana medidas \\
8 & Escritura del servicio, dashboard, Nginx, systemd e instalador; pruebas en el host & 11/11 pruebas del protocolo y JSON válido \\
9 & Intento de SSH a la Raspberry Pi con usuarios supuestos & Rechazado; sin credenciales no se pudo probar \\
10 & Reescritura de la documentación y commit local & Push no realizado: faltaban credenciales de GitHub \\
\bottomrule
\end{tabular}
\caption{Acciones del turno 2.}
\end{table}

\paragraph{Respuesta final del asistente (resumen fiel).} Firmware compilado, cargado y funcionando en el ESP32-S3 con los 4 presets (incluido 41\,667 Sa/s por canal), CRC coincidente entre chip y host y pruebas del protocolo en el host. Servicio, dashboard, Nginx y systemd escritos pero sin probar por falta de acceso SSH. Pendiente del usuario: autenticar \texttt{git push}, dar acceso a la Raspberry Pi, conectar señales conocidas a GPIO1 y GPIO2 y completar nombres y cédulas.

\subsection{Turno 3 (2 de noviembre de 2026)}

\paragraph{Usuario (literal, credenciales omitidas).}
\begin{quote}
El usuario de la raspberry es [omitido] y su contraseña es [omitido], todavia no tengo conectada la placa raspberry con la esp32, quiero que los documentos esten en LATEX y bien organizados y ademas continua con la programacion de la raspberry pi 4
\end{quote}

\paragraph{Acciones con herramientas (resumen).}

\begin{table}[H]
\centering
\small
\begin{tabular}{@{}p{0.6cm} p{7.2cm} p{7.2cm}@{}}
\toprule
\textbf{\#} & \textbf{Acción} & \textbf{Resultado} \\
\midrule
1 & Se intentó instalar una llave SSH permanente en la Raspberry Pi & \textbf{Bloqueado} por el sistema de permisos (no solicitado por el usuario); se usó la contraseña solo en cada conexión, sin guardarla \\
2 & Reconocimiento de la Raspberry Pi por SSH & Raspberry Pi 4 Model B Rev 1.2, kernel 6.18, \texttt{/dev/spidev0.0} habilitado, \texttt{gcc} disponible \\
3 & Compilación y pruebas unitarias del protocolo en la Raspberry Pi & 11/11 correctas (aarch64) \\
4 & Medición del \texttt{ioctl} SPI a 1, 8, 15{,}6, 20 y 31{,}25 MHz solicitados (sin ESP32-S3) & Modo 0 y 8 bits confirmados; tiempos del Documento 2 \\
5 & Nuevas opciones del servicio: \texttt{--bench-secs} y \texttt{--demo} (ventanas sintéticas) & Permite probar toda la cadena web sin el ESP32-S3 \\
6 & Despliegue de Nginx, sitio y servicio transitorio; pruebas HTTP, escritura atómica, refresco y reinicio & Todas correctas (Documento 2) \\
7 & Intento de descargar un compilador \LaTeX{} (Tectonic) para validar los documentos & \textbf{Bloqueado} por el sistema de permisos; los documentos no se compilaron \\
8 & Reorganización de los documentos en \LaTeX{} con preámbulo y portada compartidos & Carpeta \texttt{Documentos/Latex/} \\
\bottomrule
\end{tabular}
\caption{Acciones del turno 3.}
\end{table}

\end{document}
